% =============================================================================
% Kapitel 1: Kontext & Problemstellung
% =============================================================================

\section{Kontext und Problemstellung}

ROSI (\textit{Realtime Optical Surface Inspection}) ist ein von der Grenzebach BSH GmbH
entwickeltes Inline-Scan-System zur automatisierten Qualitätskontrolle von Baustoffen.
Im laufenden Produktionsprozess werden Materialoberflächen erfasst, Defekte per
KI-Inferenz erkannt und das Material anschließend regelbasiert in Qualitätsstufen
klassifiziert. Der Betrieb erfolgt on-premise direkt in der Produktionsanlage des
Kunden im 24/7-Dauerbetrieb.

Technisch ist ROSI als Python-Multiprocessing-System aufgebaut: Mehrere langlebige
Prozesse verarbeiten die Daten in einer Shared-Memory-basierten Pipeline, ergänzt
durch Queues und flüchtigen Zwischenspeicher (z.\,B. tmpfs/RAM-Disk). Die in dieser
Fallstudie untersuchten Fehlerklassen setzen genau an dieser Prozess- und
Speicherarchitektur an.

ROSI befindet sich in aktiver Entwicklung; die erste Kundeninbetriebnahme ist für
Mitte bis Ende 2026 geplant. Mit dem bevorstehenden Produktiveinsatz steigen die
Anforderungen an Verfügbarkeit und Fehlertoleranz erheblich: Ausfälle können
Produktionsstillstände verursachen, Inspektionsergebnisse ausbleiben lassen und im
schlimmsten Fall Qualitätsmängel unbemerkt passieren.

\subsection{Zielsetzung}

Ziel dieser Fallstudie ist es, sicherheitskritische Fehlerszenarien im industriellen
Dauerbetrieb auf Prozess- und Betriebssystemebene zu identifizieren, reproduzierbar
zu simulieren und anhand messbarer Kriterien systematisch zu bewerten.

Im Fokus stehen stille, schleichende und kaskadierende Fehler, die im
Entwicklungsbetrieb selten sichtbar werden, im 24/7-Dauerbetrieb jedoch
systemkritische Auswirkungen entfalten können.
